\chapter{Transition profiles and phase-field energy}\label{sec:phase-energy}
We separate the elementary energy identities from the later convergence of interfaces. The reference is \cite{ChAC}, Sections 2--3, in the June 29, 2023 revision. Throughout, $W(s)=(1-s^2)^2/4$ and $c_W=2\sqrt2/3$, as in Definition \ref{def:allen-cahn-energy}. The Laplacian has the sign $\Delta=\operatorname{div}\nabla$.

\begin{definition}[Local diffuse-interface energy]\label{def:phase-local-energy}
For a bounded open set $U\subset\mathbb R^n$, $n\geq1$, and $u\in H^1(U)\cap L^4(U)$, define the energy measure and its mass by
\[
 d\mu_\varepsilon^u=\left(\frac\varepsilon2|Du|^2+\frac{W(u)}\varepsilon\right)dx,
 \qquad E_\varepsilon(u;A)=\mu_\varepsilon^u(A)
\]
for Borel $A\subset U$. As a functional on $L^1(U)$, $E_\varepsilon(u;U)=+\infty$ outside this Sobolev domain. No boundary value is prescribed by this definition. This is the localized form of \cite{ChAC}, Definition 2.1.
\end{definition}

\begin{proposition}[Truncation to the two wells decreases energy]\label{prop:phase-truncation}
Let $T(s)=\max\{-1,\min\{s,1\}\}$. For finite-energy $u$ on $U$,
\[
 E_\varepsilon(T(u);A)\leq E_\varepsilon(u;A)
\]
for every Borel $A\subset U$. If a prescribed Sobolev trace already takes values in $[-1,1]$, truncation preserves that trace.
\end{proposition}
\begin{proof}
The Sobolev chain rule gives $|D(T(u))|\leq|Du|$ almost everywhere. The potential is unchanged where $|u|\leq1$, and becomes zero where $|u|>1$. Thus both energy terms decrease pointwise. Trace compatibility with Lipschitz composition gives the final assertion. This is a useful variational consequence of the potential chosen in \cite{ChAC}, Definition 2.1; it does not require the Euler equation.
\end{proof}

\begin{proposition}[Maximum bound for smooth critical points]\label{prop:phase-maximum-bound}
A smooth solution of $\varepsilon^2\Delta_g u=u(u^2-1)$ on a closed Riemannian manifold satisfies $-1\leq u\leq1$. The same conclusion holds for $u\in C^2(U)\cap C^0(\overline U)$ on a bounded Euclidean domain when its boundary values lie in $[-1,1]$.
\end{proposition}
\begin{proof}
A maximum larger than one must occur at an interior point, where $\Delta_g u\leq0$, while $u(u^2-1)>0$. This contradicts the equation. At a minimum below minus one the two signs are reversed, giving the other bound. Compactness supplies the extrema; connectedness is unnecessary. This proves \cite{ChAC}, Exercise 2.1(b), and the stated boundary version.
\end{proof}

\begin{proposition}[Zero energy and connected components]\label{prop:phase-zero-energy}
On a connected bounded open Euclidean domain, a finite-energy function has zero energy if and only if it equals $+1$ or $-1$ almost everywhere. On a closed manifold the same assertion holds on each connected component; different components may choose different signs.
\end{proposition}
\begin{proof}
Both integrands are nonnegative. Zero energy implies $Du=0$ and $W(u)=0$ almost everywhere. A Sobolev function with zero weak gradient is constant on each connected component: mollification makes it constant on interior balls, and overlapping balls propagate the constant. The zeros of $W$ are precisely $\pm1$. Conversely either constant has zero energy. The argument in coordinate charts proves the manifold statement. The connectedness qualification refines the observation following Exercise 2.1 in \cite{ChAC}.
\end{proof}

\begin{lemma}[The one-dimensional Hamiltonian]\label{lem:phase-hamiltonian}
If $u\in C^2(I)$ solves $\varepsilon^2u''=W'(u)$ on an interval, then
\[
 \frac{\varepsilon^2}{2}(u')^2-W(u)
\]
is constant on $I$.
\end{lemma}
\begin{proof}
Its derivative is $u'(\varepsilon^2u''-W'(u))=0$. The identity is the conservation law in \cite{ChAC}, Section 2.2. It applies to solutions of the ODE, whereas the energy inequalities below apply to arbitrary competitors.
\end{proof}

\begin{theorem}[Classification of finite-energy entire transition profiles]\label{thm:phase-ode-classification}
Let $u\in C^2(\mathbb R)$ solve $\varepsilon^2u''=W'(u)$. Its total one-dimensional energy is finite if and only if it is one of
\[
 u\equiv\pm1,\qquad
 u(t)=\pm\tanh\left(\frac{t-t_0}{\varepsilon\sqrt2}\right),
 \quad t_0\in\mathbb R.
\]
Every nonconstant profile in this list has energy $c_W$.
\end{theorem}
\begin{proof}
Finiteness gives points $t_j\to\infty$ at which both $(u'(t_j))^2$ and $W(u(t_j))$ tend to zero. Lemma \ref{lem:phase-hamiltonian} therefore has constant zero, so $\varepsilon|u'|=\sqrt{2W(u)}$. If $u$ reaches either well, its derivative there is zero; uniqueness for the smooth second-order ODE makes it constant. Otherwise it stays in one of $(-\infty,-1)$, $(-1,1)$, or $(1,\infty)$ and $u'$ never vanishes.

In the middle interval, separation of variables gives exactly the two translated hyperbolic-tangent families. In either exterior interval, choose the time direction in which $|u|$ increases. The time to reach infinity is finite, because $\int_a^\infty(s^2-1)^{-1}ds<\infty$ for $a>1$. This contradicts an entire solution. Finally the equipartition identity and substitution $s=u(t)$ give energy $\int_{-1}^1\sqrt{2W(s)}ds=c_W$ for the nonconstant profiles; the constants have zero energy. This supplies the details of \cite{ChAC}, Lemma 2.3.
\end{proof}

\begin{definition}[The phase primitive]\label{def:phase-primitive}
Set
\[
 \Phi(t)=\int_0^t\sqrt{2W(s)}\,ds
        =\frac1{\sqrt2}\int_0^t|1-s^2|\,ds.
\]
It is odd, continuously differentiable, strictly increasing, and maps $\mathbb R$ onto $\mathbb R$. Its inverse is continuous. In particular $\Phi(\pm1)=\pm c_W/2$ and $|\Phi(t)|\leq C(1+|t|^3)$. Strict increase follows because its derivative vanishes only at two points. This is the primitive in \cite{ChAC}, Section 3.1.
\end{definition}

\begin{lemma}[The Sobolev primitive and Modica's energy inequality]\label{lem:phase-modica-energy}
If $U$ is bounded and $u\in H^1(U)\cap L^4(U)$, then $\Phi(u)\in W^{1,1}(U)$ and
\[
 D\Phi(u)=\sqrt{2W(u)}Du,\qquad
 |D\Phi(u)|\leq\frac\varepsilon2|Du|^2+\frac{W(u)}\varepsilon
 \quad\text{almost everywhere}.
\]
Consequently $|D\Phi(u)|(A)\leq E_\varepsilon(u;A)$ for every Borel $A\subset U$.
\end{lemma}
\begin{proof}
The cubic growth gives $\Phi(u)\in L^1$. Truncate $u$ to $[-k,k]$ and apply the Lipschitz Sobolev chain rule to the corresponding bounded-range primitive. The gradients converge in $L^1$, since they are bounded by $C(1+|u|^2)|Du|$, which is integrable by Cauchy--Schwarz and $u\in L^4$. The functions converge in $L^1$ by the cubic bound. Passing to weak derivatives proves the chain rule. Finally apply $2ab\leq a^2+b^2$ with $a=\sqrt{\varepsilon/2}|Du|$ and $b=\sqrt{W(u)/\varepsilon}$. This proves \cite{ChAC}, equation (3.1), including its Sobolev justification. It is an energy inequality, not the pointwise gradient estimate for entire solutions also associated with Modica's name.
\end{proof}

\begin{proposition}[A sharp bound for a one-dimensional transition]\label{prop:phase-transition-cost}
For $u\in H^1(a,b)$, using its continuous representative,
\[
 \int_a^b\left(\frac\varepsilon2|u'|^2+\frac{W(u)}\varepsilon\right)dt
 \geq |\Phi(u(b))-\Phi(u(a))|.
\]
Equality holds exactly when $\varepsilon|u'|=\sqrt{2W(u)}$ almost everywhere and $\Phi(u)$ is monotone. For an entire finite-energy function with limits $-1$ and $+1$, the analogous bound is $c_W$ and is attained by the increasing profile in Theorem \ref{thm:phase-ode-classification}.
\end{proposition}
\begin{proof}
Integrate Lemma \ref{lem:phase-modica-energy}, then use the fundamental theorem for the absolutely continuous function $\Phi(u)$. Equality in the first inequality requires $a=b$ in the square inequality, and equality between total variation and endpoint difference requires a derivative of constant sign almost everywhere. These are exactly the stated conditions. On the whole line, apply the bound on $(-R,R)$ and let $R\to\infty$. This makes the transition cost from \cite{ChAC}, Exercise 2.3, explicit for competitors as well as solutions.
\end{proof}

\begin{definition}[Energy discrepancy]\label{def:phase-discrepancy}
For a finite-energy function on $U$, its discrepancy is the signed measure
\[
 d\xi_\varepsilon^u=
 \left(\frac\varepsilon2|Du|^2-\frac{W(u)}\varepsilon\right)dx.
\]
It compares gradient and potential energy. Vanishing is equipartition; no sign is asserted for a general function. This notation isolates the two terms appearing in \cite{ChAC}, Sections 2.2 and 3.1.
\end{definition}

\begin{proposition}[The energy defect controls discrepancy]\label{prop:phase-defect}
Let $E=E_\varepsilon(u;U)$ and $D=E-\int_U|D\Phi(u)|$. Then
\[
 D=\int_U\left(\sqrt{\frac\varepsilon2}|Du|-\sqrt{\frac{W(u)}\varepsilon}\right)^2dx,
 \qquad |\xi_\varepsilon^u|(U)\leq\sqrt{2ED}.
\]
In particular $D\geq0$, and $D=0$ is equivalent to equality of the two energy densities almost everywhere.
\end{proposition}
\begin{proof}
With $a,b$ as in Lemma \ref{lem:phase-modica-energy}, the defect is $\int(a-b)^2$. The discrepancy density is $(a-b)(a+b)$. Cauchy--Schwarz and $(a+b)^2\leq2(a^2+b^2)$ prove the bound. This quantitative refinement of the square inequality will turn convergence of total energy into equipartition.
\end{proof}
